\documentclass{article}

\usepackage[british]{babel}
\usepackage[T1]{fontenc}
\usepackage{lmodern}
\usepackage{microtype}

\usepackage{amsmath}
\usepackage{amssymb}
\usepackage{amsthm}

\usepackage{fullpage}
\usepackage{parskip}
\usepackage{fancyhdr}
\usepackage{enumitem}

\usepackage{xcolor}
\usepackage{hyperref}
\usepackage[capitalise,nameinlink]{cleveref}

\hypersetup{
  colorlinks=true,
  allcolors=blue!50!black,
  pdftitle={The Sylow argument},
  pdfauthor={Joe Lamond and Joshua Munchau},
}

\pagestyle{fancy}
\fancyhf{}
\fancyhead[L]{The Sylow argument}
\fancyhead[R]{Joe Lamond and Joshua Munchau}
\fancyfoot[C]{\thepage}
\setlength{\headheight}{14pt}
\setlength{\headsep}{10pt}

\newtheorem{theorem}{Theorem}
\newtheorem*{remark}{Remark}

\DeclareMathOperator{\Gal}{Gal}
\DeclareMathOperator{\Aut}{Aut}

\begin{document}

\section*{The Sylow argument}
This is Step~3 of the proof sketch in the \emph{Outline of formalisation}.
\begin{theorem}\label{thm:the-sylow-argument}
Suppose $R$ is a real closed field, and $R\subseteq L$ is a finite Galois extension. Then, $L$ has degree a power of $2$ over $R$.
\end{theorem}
The proof relies on the fact that every finite extension of odd degree of $R$
has degree 1.
\begin{proof}
Let $G=\Aut_R(L)$. Write $|G|=2^nm$, where $m$ is odd. By the Sylow Theorems,
$G$ has a subgroup $P$ of order $2^n$. Let $E=L^P$. The Galois correspondence
tells us that $P=\Aut_E(L)$. Since $R\subseteq L$ and $E\subseteq L$ are both
finite Galois, we have $[L:R]=|G|=2^nm$ and $[L:E]=|P|=2^n$. The tower law then
implies that $[E:R]=m$ is odd. Hence $[E:R]=1$, completing the proof.
\end{proof}

\end{document}